\textbf{Phần 1. Trắc nghiệm lựa chọn}
\setcounterpageref{ex}{0}

\begin{ex} %[2D14V68]
Số liền trước của số lớn nhất có ba chữ số là
\choice
{$999$}
{$1 000$}
{\True $998$}
{$997$}
\loigiai{
Số lớn nhất có ba chữ số là 999. Số liền trước là $999 - 1 = 998$.
\textbf{Đáp án của câu hỏi là C}
}
\end{ex}

\begin{ex} %[2D14V68]
Số gồm $7$ trăm, $0$ chục và $4$ đơn vị viết là
\choice
{\True $704$}
{$740$}
{$74$}
{$407$}
\loigiai{
Viết là $704$.
\textbf{Đáp án của câu hỏi là A}
}
\end{ex}

\begin{ex} %[2D14V68]
Chữ số $5$ trong số $358$ có giá trị bằng bao nhiêu?
\choice
{$500$}
{$5$}
{$15$}
{\True $50$}
\loigiai{
Chữ số 5 nằm ở hàng chục nên có giá trị là 50.
\textbf{Đáp án của câu hỏi là D}
}
\end{ex}

\begin{ex} %[2D14V68]
Từ các chữ số $3, 0, 8$, số chẵn lớn nhất có ba chữ số khác nhau lập được là
\choice
{$803$}
{\True $830$}
{$380$}
{$308$}
\loigiai{
Số lớn nhất có chữ số hàng trăm lớn nhất (8), chục (3), đơn vị chẵn (0): số 830.
\textbf{Đáp án của câu hỏi là B}
}
\end{ex}

\begin{ex} %[2D14V68]
Số tròn trăm liền sau của số $450$ là
\choice
{$400$}
{$460$}
{\True $500$}
{$600$}
\loigiai{
Số tròn trăm lớn hơn và liền kề sau 450 là 500.
\textbf{Đáp án của câu hỏi là C}
}
\end{ex}

\begin{ex} %[2D14V68]
Cho $x$ là số có ba chữ số sao cho $789 < x < 792$. Có bao nhiêu giá trị của $x$?
\choice
{\True $2$ số}
{$3$ số}
{$1$ số}
{$4$ số}
\loigiai{
$x \in \{790, 791\}$, có đúng 2 số.
\textbf{Đáp án của câu hỏi là A}
}
\end{ex}

\begin{ex} %[2D14V68]
Tổng của chữ số hàng trăm và chữ số hàng chục của số $625$ là
\choice
{$7$}
{$11$}
{$13$}
{\True $8$}
\loigiai{
$6 + 2 = 8$.
\textbf{Đáp án của câu hỏi là D}
}
\end{ex}

\begin{ex} %[2D14V68]
Điền dấu thích hợp: $500 + 40 + 9 \dots 500 + 50$.
\choice
{$>$}
{\True $<$}
{$=$}
{Không so sánh được}
\loigiai{
$500 + 40 + 9 = 549 < 550 = 500 + 50$.
\textbf{Đáp án của câu hỏi là B}
}
\end{ex}

\begin{ex} %[2D14V68]
Tìm số tự nhiên có ba chữ số, biết chữ số hàng trăm là số lớn nhất có một chữ số, hàng chục kém hàng trăm 3 đơn vị và hàng đơn vị bằng 0.
\choice
{$930$}
{$850$}
{\True $960$}
{$970$}
\loigiai{
Hàng trăm là 9, hàng chục là $9 - 3 = 6$, hàng đơn vị là 0: số 960.
\textbf{Đáp án của câu hỏi là C}
}
\end{ex}

\begin{ex} %[2D14V68]
Có bao nhiêu số tròn chục có ba chữ số có chữ số hàng trăm bằng $2$?
\choice
{\True $10$ số}
{$9$ số}
{$8$ số}
{$11$ số}
\loigiai{
Các số là: 200, 210, 220, ..., 290. Có tất cả 10 số.
\textbf{Đáp án của câu hỏi là A}
}
\end{ex}

\begin{ex} %[2D14V68]
Khi đổi chỗ chữ số hàng trăm và chữ số hàng đơn vị của số $437$, ta được số mới là
\choice
{$743$}
{$374$}
{$473$}
{\True $734$}
\loigiai{
Đổi vị trí 4 và 7 ta được 734.
\textbf{Đáp án của câu hỏi là D}
}
\end{ex}

\begin{ex} %[2D14V68]
Hiệu giữa số bé nhất có bốn chữ số và số lớn nhất có hai chữ số là
\choice
{$900$}
{\True $901$}
{$910$}
{$899$}
\loigiai{
$1 000 - 99 = 901$.
\textbf{Đáp án của câu hỏi là B}
}
\end{ex}

\textbf{Phần 2. Trắc nghiệm đúng sai}
\setcounterpageref{ex}{0}

\begin{ex} %[2D14V68]
Xét các khẳng định về số học phạm vi 1 000:
\choiceTF
{\True Số 1 000 gồm 10 trăm hoặc 100 chục.}
{Số 305 đọc là ba trăm năm mươi.}
{\True Dãy số 120, 140, 160, 180 là dãy các số chẵn tăng dần.}
{\True Tổng các chữ số của số 456 bằng 15.}
\loigiai{
\begin{itemize}
    \item Ý a) Đúng: Đúng.
    \item Ý b) Sai: 305 đọc là ba trăm linh năm.
    \item Ý c) Đúng: Đúng.
    \item Ý d) Đúng: 4 + 5 + 6 = 15. Đúng.
\end{itemize}
\begin{itemize}
    \item a) \textbf{Đúng}.
    \item b) \textbf{Sai}.
    \item c) \textbf{Đúng}.
    \item d) \textbf{Đúng}.
\end{itemize}
}
\end{ex}

\begin{ex} %[2D14V68]
Cho ba số tự nhiên A = 345, B = 354, C = 435:
\choiceTF
{\True A < B < C.}
{\True Số lớn nhất trong ba số là C.}
{Hiệu giữa B và A bằng 11.}
{\True Cả ba số đều có tổng các chữ số bằng 12.}
\loigiai{
\begin{itemize}
    \item Ý a) Đúng: 345 < 354 < 435. Đúng.
    \item Ý b) Đúng: 435 là lớn nhất. Đúng.
    \item Ý c) Sai: 354 - 345 = 9, không phải 11. Sai.
    \item Ý d) Đúng: 3 + 4 + 5 = 12. Đúng.
\end{itemize}
\begin{itemize}
    \item a) \textbf{Đúng}.
    \item b) \textbf{Đúng}.
    \item c) \textbf{Sai}.
    \item d) \textbf{Đúng}.
\end{itemize}
}
\end{ex}

\textbf{Phần 3. Trả lời ngắn}
\setcounterpageref{ex}{0}

\begin{ex} %[2D14V68]
Tìm số tự nhiên có ba chữ số khác nhau bé nhất.
\shortans{102}
\loigiai{
\textbf{Đáp số:} $102$\\
Hàng trăm chọn 1, chục chọn 0, đơn vị chọn 2: số 102.
}
\end{ex}

\begin{ex} %[2D14V68]
Tìm số liền trước của số tròn trăm bé nhất có ba chữ số.
\shortans{99}
\loigiai{
\textbf{Đáp số:} $99$\\
Số tròn trăm bé nhất có 3 chữ số là 100. Liền trước là 99.
}
\end{ex}

\begin{ex} %[2D14V68]
Có bao nhiêu số có ba chữ số có tận cùng là chữ số 5 và chữ số hàng trăm là 8?
\shortans{10}
\loigiai{
\textbf{Đáp số:} $10$\\
Dạng $8a5$ với $a \in \{0, 1, \dots, 9\}$. Có 10 số.
}
\end{ex}

\begin{ex} %[2D14V68]
Tìm một số biết rằng lấy số đó cộng với 250 thì được số tròn trăm lớn nhất có ba chữ số.
\shortans{650}
\loigiai{
\textbf{Đáp số:} $650$\\
Số tròn trăm lớn nhất có 3 chữ số là 900. $900 - 250 = 650$.
}
\end{ex}

\begin{ex} %[2D14V68]
Viết số gồm 5 trăm, 8 chục và 3 đơn vị.
\shortans{583}
\loigiai{
\textbf{Đáp số:} $583$\\
Viết là 583.
}
\end{ex}

\begin{ex} %[2D14V68]
Tìm hiệu của số lớn nhất có ba chữ số và số bé nhất có ba chữ số.
\shortans{899}
\loigiai{
\textbf{Đáp số:} $899$\\
$999 - 100 = 899$.
}
\end{ex}
